\section{Related Work}

\paragraph{Efficient CoT} While Chain-of-Thought (CoT) enhances the reasoning performance of LLMs, it introduces significant computational overhead. Researchers have sought methods to reduce this overhead while retaining the benefits of CoT. One intuitive approach is to simplify~\cite{marconato2024not}, skip~\cite{Ding:2024cotshortcut, liu2024skipstep}, or generate reasoning steps in parallel~\cite{ning2023skeleton}. Another research direction involves compressing CoTs into latent representations~\cite{goyal2024think, deng2024explicit, hao2024training, cheng2024compressed}, allowing LLMs to reason without explicitly generating discrete tokens. To mitigate CoT redundancy, \citet{han2024token} guides token consumption through dynamic token budget estimation. \citet{kang2024c3ot} prompts GPT-4 to shorten CoT trajectories, and then fine-tunes LLMs using compressed CoTs. In contrast, this work focuses on pruning CoT tokens based on their semantic importance. Moreover, \method leverages a small LM for token pruning, significantly reducing computational overhead.

\paragraph{Prompt Compression} The growing demand for long-context prompts has led to substantial computational and memory challenges. To address this, researchers have explored various prompt compression techniques. One intuitive approach involves using a lightweight LM to generate more concise prompts~\cite{chuang-etal-2024-learning}. Considering that natural language formats inevitably contain redundancy, some studies have introduced implicit continuous tokens to represent long-context inputs~\cite{chevalier-etal-2023-adapting, ge2024incontext, mohtashami2023randomaccess}. Another line of research focuses on directly compressing prompts by filtering low-informative tokens~\cite{li:2023selective, jiang2023:llmlingua, pan:2024llmlingua2}. For instance, Selective Context uses the perplexity of LLMs to measure token importance and removes less important tokens. LLMLingua-2~\cite{pan:2024llmlingua2} introduces a small bidirectional language model for token importance measurement and trains this LM with \texttt{GPT-4} compression data, which serves as the token importance metric in this work.